\documentclass{ctexbook}
\usepackage{Note}
\usepackage{unicode-math}
\setCJKmainfont[
  BoldFont = 思源宋体 Bold,
  ItalicFont = 楷体,
]{宋体}
\setcounter{secnumdepth}{3}
\begin{document}
\chapter{原子结构与光谱}
\section{类氢原子}
\subsection{类氢原子的薛定谔方程}
氢原子的哈密顿量包含三部分:
\[\hat{H}=\underbrace{-\dfrac{\hbar^2}{2m_e}\nabla_{\vec{r}_e}^2}_{\hat{E}_{k,\text{electron}}}\underbrace{-\dfrac{\hbar^2}{2m_N}\nabla_{\vec{r}_N}^2}_{\hat{E}_{k,\text{nucleus}}}\underbrace{-\dfrac{Ze^2}{4\pi\ep_0 r}}_{\hat{V}}\]
其中$\vec{r}_e$, $\vec{r}_N$分别是电子和原子核的坐标, $Z$是核电荷数, $r$是电子与原子核的距离.为了求解该方程,改写动能项为质心动能与相对动能之和:
\[-\dfrac{\hbar^2}{2m_e}\nabla_{\vec{r}_e}^2-\dfrac{\hbar^2}{2m_N}\nabla_{\vec{r}_N}^2=-\dfrac{\hbar^2}{2M}\nabla_{\vec{R}}^2-\dfrac{\hbar^2}{2\mu}\nabla_{\vec{r}}^2\]
其中
\[M=m_e+m_N,\quad\mu=\dfrac{m_em_N}{m_e+m_N},\quad\vec{R}=\dfrac{m_e\vec{r}_e+m_N\vec{r}_N}{m_e+m_N},\quad\vec{r}=\vec{r}_e-\vec{r}_N\]
这样体系的定态薛定谔方程即可改写为
\[-\dfrac{\hbar^2}{2M}\nabla_{\vec{R}}^2\iPsi-\dfrac{\hbar^2}{2\mu}\nabla_{\vec{r}}^2\iPsi-\dfrac{Ze^2}{4\pi\ep_0r}\iPsi=E_{\text{tot}}\iPsi\]
分量变量,令
\[\iPsi(\vec{r},\vec{R})=\chi(\vec{R})\psi(\vec{r}),\quad E_{\text{tot}}=\mathcal{E}+E\]
则可以得到
\[\left\{\begin{array}{l}
    -\dfrac{\hbar^2}{2M}\nabla_{\vec{R}}^2\chi=\mathcal{E}\chi\\[6pt]
    -\dfrac{\hbar^2}{2\mu}\nabla_{\vec{r}}^2\psi-\dfrac{Ze^2}{4\pi\ep_0 r}\psi=E\psi
\end{array}\right.\]
\indent 质心波函数$\chi$满足自由粒子运动方程,其解为
\[\chi(\vec{R})=\dfrac{1}{(2\pi)^{3/2}}\e^{-\i\vec{k}\cdot\vec{R}},\quad\mathcal{E}=\dfrac{\hbar^2k^2}{2M}\]
下面求解相对运动方程.
\subsection{相对运动方程}
在相对运动方程中,由于$m_N\gg m_e$,因此$\mu\approx m_e$.利用球坐标中拉普拉斯算子的表达形式
\[\nabla^2=\dfrac{1}{r}\dfrac{\p^2}{\p r^2}r+\dfrac{1}{r^2}\hat{\Lambda}^2,\quad\hat{\Lambda}^2=\dfrac{1}{\sin^2\theta}\dfrac{\p}{\p\phi^2}+\dfrac{1}{\sin\theta}\dfrac{\p}{\p\theta}\sin\theta\dfrac{\p}{\p\theta}\]
对波函数$\psi(\vec{r})$分离变量
\[\psi(\vec{r})=\psi(r,\theta,\phi)=R(r)Y(\theta,\phi)\]
可得
\[-\dfrac{\hbar^2}{2\mu}\left[\dfrac1r\left(\dfrac{\p^2}{\p r^2}\right)+\dfrac{\hat{\Lambda}^2}{r^2}\right]R(r)Y(\theta,\phi)-\dfrac{Ze^2}{4\pi\ep_0 r}R(r)Y(\theta,\phi)=ER(r)Y(\theta,\phi)\]
对上式两端同除以$R(r)Y(\theta,\phi)$可得
\[-\dfrac{\hbar^2}{2\mu}\dfrac{1}{rR(r)}\left(\dfrac{\p^2}{\p r^2}rR(r)\right)-\dfrac{\hbar^2}{2\mu}\dfrac{1}{r^2Y(\theta,\phi)}\hat{\Lambda}^2Y(\theta,\phi)-\dfrac{Ze}{4\pi\ep_0r}=E\]
对上式两端同乘以$r^2$可得
\[-\dfrac{\hbar^2}{2\mu}\dfrac{r}{R(r)}\left(\dfrac{\p^2}{\p r^2}rR(r)\right)-\dfrac{Zer}{4\pi\ep_0}-Er^2=\dfrac{\hbar^2}{2\mu}\dfrac{1}{Y(\theta,\phi)}\hat{\Lambda}^2Y(\theta,\phi)=\text{Const}\]
为了方便求解(主要是角方程与角动量量子化的统一),令常数为$C\dfrac{\hbar^2}{2\mu}$.
\subsubsection{角向方程}
于是角向方程为
\[\hat{\Lambda}^2Y(\theta,\phi)=CY(\theta,\phi)\]
我们已经知道$\hat{\Lambda}^2$的本征函数是球谐函数$Y_{l,m_l}$,本征值为$-l(l+1)$.于是角向方程即得其解.
\subsubsection{径向方程}
在$C=-l(l+1)$的前提下,径向方程可以改写为
\[-\dfrac{\hbar^2}{2\mu}\left(\dfrac{\p^2 R}{\p r^2}+\dfrac{2}{r}\dfrac{\p R}{\p r}\right)+\left[\dfrac{l(l+1)\hbar^2}{2\mu r^2}-\dfrac{Ze^2}{4\pi\ep_0r}\right]R=ER\]
做变量代换
\[\rho=\dfrac{2Zr}{na},\quad a=\dfrac{m_e}{\mu}a_0,\quad a_0=\dfrac{4\pi\ep_0\hbar^2}{m_ee^2},\quad R(r)=r^l L(r)\e^{-\alpha r}\]
则本征波函数满足
\[R_{n,l}(r)=N_{n,l}\rho^l L_{n,l}(\rho)\e^{-\rho/2}\]
其中$L_{n,l}$时连带拉盖尔多项式,满足
\[x\dfrac{\d^2 L_{n,l}}{\d x^2}+(l+1-x)\dfrac{\d L_{n,l}}{\d x}+nL=0\]
$n$必须是整数,称作\tbf{主量子数},并且$l=0,1,\cdots,n-1$.径向波函数$R(r)$除去$r=0$之外的零点称作\tbf{节点},节点的数目为$n-l-1$.
\subsection{原子单位制}
考虑类氢原子的薛定谔方程(忽略原子核动能项):
\[-\dfrac{\hbar^2}{2m_e}\nabla^2\psi-\dfrac{Ze^2}{4\pi\ep_0 r}\psi=E\psi\]
于是
\[-\dfrac12\nabla^2\psi-\dfrac{Zm_e e^2}{4\pi\ep_0\hbar^2 r}=\dfrac{m_e}{\hbar^2}E\psi\]
令$a_0=\dfrac{4\pi\ep_0\hbar^2}{m_e e^2}$为Bohr半径,则
\[-\dfrac12\nabla^2\psi-\dfrac{Z}{a_0 r}\psi=\dfrac{m_e}{\hbar^2}E\psi\]
令$\tilde{\vec{r}}=\dfrac{\vec{r}}{a_0}$,则有
\[-\dfrac{1}{2}\nabla_{\tilde{\vec{r}}}^2\psi-\dfrac{Z}{\tilde{r}}\psi=\dfrac{m_e a_0^2}{\hbar^2}E\psi\]
令$\tilde{E}=\dfrac{m_e a_0^2}{\hbar^2}E=\dfrac{4h^2\ep_0^2}{m_e e^4}E$,则有
\[-\dfrac12\nabla^2\psi-\dfrac{Z}{r}\psi=E\psi\]
这里省去了$\sim$符号.在原子单位中长度$a_0=1$,质量$m_e=1$, $\hbar=1$.由此,氢原子1s轨道的能量为$-\dfrac12$.
\section{多电子原子}
\subsection{自旋}
Stern-Gerlach实验表明银原子束通过非匀强磁场散射,并且偏转程度并不连续分布,只有两种模式.实验证明电子的\tbf{自旋量子数}$s=\dfrac12$,自旋角动量的模为
\[S=\hbar\sqrt{s(s+1)}\]
自旋量子数是粒子的内禀属性.\\
\indent 电子的自旋状态有两种,对应的\tbf{自旋磁量子数}分别为$m_s=\dfrac12$和$m_s=-\dfrac12$.两种状态通常用$\alpha$和$\beta$,或向上和向下的箭头标记.
\indent 根据自旋量子数的不同可以把粒子分为两类:
\begin{definition}[费米子和玻色子]
	自旋量子数为整数的粒子是\tbf{玻色子},自旋量子数为半整数的粒子是\tbf{费米子}.
\end{definition}
\subsubsection{交换对称性}
玻色子和费米子具有如下重要性质:
\begin{theorem}
	费米子的多粒子波函数满足交换反对称性,即
	\[\iPsi_F(1,2)=-\iPsi_F(2,1)\]
	玻色子的多粒子波函数满足交换对称性,即
	\[\iPsi_B(1,2)=\iPsi_B(2,1)\]
\end{theorem}
\subsection{多电子原子}
氦原子的定态薛定谔方程为
\[\hat{H}=\underbrace{-\dfrac{\hbar^2}{2m_e}\nabla_{\vec{r}_1}^2-\dfrac{2e^2}{4\pi\ep_0 r_1}}_{\hat{H}_1}\underbrace{-\dfrac{\hbar^2}{2m_e}\nabla_{\vec{r}_2}^2-\dfrac{2e^2}{4\pi\ep_0 r_2}}_{\hat{H}_2}+\dfrac{e^2}{4\pi\ep_0r_{12}}\]
如果电子之间没有相互作用,则可以对波函数分离变量$\iPsi(\vec{r}_1,\vec{r}_2)=\psi_1(\vec{r}_1)\psi_2(\vec{r}_2)$,将问题转化为两个类氢原子问题.这种方法就是\tbf{单电子近似}或\tbf{轨道近似}.\\
\indent 除去空间部分的波函数,还需要考虑电子的自旋,以及作为费米子的交换反对称性.将氦原子的全波函数的空间部分与自旋部分进行变量分离:
\[\iPsi(1,2)=\psi(\vec{r}_1,\vec{r}_2)\chi(1,2)\]
根据泡利不相容原理,由于电子是费米子,因此全波函数必须是反对称的,即$\iPsi(1,2)=-\iPsi(2,1)$.在单电子近似下,基态氦原子的空间波函数为
\[\psi(\vec{r}_1,\vec{r}_2)=\psi_{1\text{s}}(\vec{r}_1)\psi_{1\text{s}}(\vec{r}_2)\]
\subsection{构造原理}
\subsubsection{钻穿与屏蔽}
多电子原子中的电子受到来自其它所有电子的库仑排斥作用.如果电子离核的距离为$r$,它受到的平均排斥作用等效于与核的位置相同的负点电荷,其电荷量为半径为$r$内所有电子的总电荷.
\subsubsection{洪特规则}
\begin{theorem}[洪特规则]
	基态原子采用具有最多未成对电子数目的组态.
\end{theorem}
考虑电子$1$, $2$处于不同轨道$a$, $b$上的波函数
\[\iPsi_{\pm}=\dfrac{1}{\sqrt2}\left[\psi_a(\vec{r}_1)\psi_b(\vec{r}_2)\pm\psi_a(\vec{r}_1)\psi_b(\vec{r}_2)\right]\]
对于对称的组合$\iPsi_+$,为了满足泡利原理,自旋波函数$\chi(1,2)$应当是反对称的,这只能有
\[\chi(1,2)=\sigma_-(1,2)=\dfrac{1}{\sqrt{2}}\left[\alpha(1)\beta(2)-\alpha(2)\beta(1)\right]\]
而对于反对称的组合$\iPsi_-$,自旋波函数$\chi(1,2)$则应当是对称的,也即
\[\alpha(1)\alpha(2),\quad\beta(1)\beta(2),\quad\sigma_+(1,2)=\dfrac{1}{\sqrt{2}}\left[\alpha(1)\beta(2)+\alpha(2)\beta(1)\right]\]
中的一者.\\
\indent 简单而言,当$\vec{r}_1\to\vec{r}_2$时,由于库伦势能增大,波函数$\iPsi(1,2)$应当趋近于$0$.反对称的组合$\iPsi_-(1,2)$天然地满足这一点,而$\iPsi_+(1,2)$则不一定.因而这两个空间波函数对应于不同的能量,我们称之为\tbf{三重态}(对应于$\iPsi_-$)和\tbf{单重态}(对应于$\iPsi_+$).
\section{原子光谱}
原子中的电子分布从一个态向另一个态转变并导致能量变化(通常是吸收/放出一个光子)的过程称作\tbf{跃迁}.在光谱学中,跃迁被描述为在两个\tbf{谱项}间发生.\\
\indent 并非所有谱项间的跃迁都能发生.如果它们能发生,则称这一跃迁是\tbf{允许}的;如果它们不发生,则光谱跃迁是\tbf{禁阻}的.\\
\indent 首先,光子的自旋角动量对应的量子数为$s=1$.为了保证自旋角动量守恒,我们引入以下近似对体系进行处理:
\begin{theorem}[偶极近似]
	如果原子(或分子)需要吸收(或产生)频率为$\nu$的光子,那么它们必须拥有在此频率上振荡的偶极子.
\end{theorem}
\subsection{微扰理论}
研究两个态之间的变化必须引入时间参量,也就是求解含时薛定谔方程.这是很难求解的,因此我们需要采用合适的办法进行近似.\tbf{微扰理论}就是一个很合适的办法.
\subsubsection{定态微扰}
\indent 首先考虑定态微扰问题:
\[\hat{H}\psi_n=E_n\psi_n,\quad\hat{H}=\hat{H}^{(0)}+\hat{H}^{(1)}\]
其中$\hat{H}^{(0)}$是对体系微扰前的哈密顿算符, $\hat{H}^{(1)}$是微扰哈密顿算符.微扰前的薛定谔方程
\[\hat{H}^{(0)}\psi_n^{(0)}=E_n^{(0)}\psi_n^{(0)}\]
的解已知.现在,对$\hat{H}$, $E_n$和$\psi_n$泰勒展开:
\[\hat{H}=\hat{H}^{(0)}+\lambda\hat{H}^{(1)}+o(\lambda)\]
\[E_n=E_n^{(0)}+\lambda E_n^{(1)}+o(\lambda)\]
\[\psi_n=\psi_n^{(0)}+\lambda\psi_n^{(1)}+o(\lambda)\]
代入薛定谔方程,展开后忽略高阶无穷小项并比较$\lambda$项的系数可得
\[\hat{H}^{(1)}\psi_n^{(0)}+\hat{H}^{(0)}\psi_n^{(1)}=E_n^{(1)}\psi_n^{(0)}+E_n^{(0)}\psi_n^{(1)}\]
\indent 将上述方程与$\psi_{n}^{(0)}$做内积可得
\[\bra{\psi_{n}^{(0)}}\hat{H}^{(1)}\ket{\psi_{n}^{(0)}}+\bra{\psi_{n}^{(0)}}\hat{H}^{(0)}\ket{\psi_{n}^{(1)}}=E_n^{(0)}\inprod{\psi_{n}^{(0)}}{\psi_{n}^{(1)}}+E_n^{(1)}\inprod{\psi_{n}^{(0)}}{\psi_{n}^{(0)}}\]
由于$\hat{H}$是厄米算符,因此$\hat{H}^{(k)}$也是厄米算符,从而
\[\bra{\psi_{n}^{(0)}}\hat{H}^{(0)}\ket{\psi_{n}^{(1)}}=\hat{H}^{(0)}\inprod{\psi_{n}^{(0)}}{\psi_{n}^{(1)}}=E_n^{(0)}\inprod{\psi_{n}^{(0)}}{\psi_{n}^{(1)}}\]
于是
\[E_n^{(1)}=\bra{\psi_{n}^{(0)}}\hat{H}^{(1)}\ket{\psi_{n}^{(0)}}\]
\indent 将上述方程与$\psi_k^{(0)}(k\neq n)$做内积可得
\[\bra{\psi_{k}^{(0)}}\hat{H}^{(1)}\ket{\psi_{n}^{(0)}}+\bra{\psi_{k}^{(0)}}\hat{H}^{(0)}\ket{\psi_{n}^{(1)}}=E_n^{(0)}\inprod{\psi_{k}^{(0)}}{\psi_{n}^{(1)}}+E_n^{(1)}\inprod{\psi_{k}^{(0)}}{\psi_{n}^{(0)}}\]
根据波函数的正交性可得$\inprod{\psi_{k}^{(0)}}{\psi_{n}^{(0)}}=0$,而
\[\bra{\psi_{k}^{(0)}}\hat{H}^{(0)}\ket{\psi_{n}^{(1)}}=\hat{H}^{(0)}\inprod{\psi_{k}^{(0)}}{\psi_{n}^{(1)}}=E_k^{(0)}\inprod{\psi_{k}^{(0)}}{\psi_{n}^{(1)}}\]
从而
\[\inprod{\psi_{k}^{(0)}}{\psi_{n}^{(1)}}=\dfrac{\bra{\psi_{k}^{(0)}}\hat{H}^{(1)}\ket{\psi_{n}^{(0)}}}{E_n^{(0)}-E_k^{(0)}}\]
于是利用基组展开的技巧可得
\[\ket{\psi_n^{(1)}}=\sum_{k}\ket{\psi_k^{(0)}}\inprod{\psi_{k}^{(0)}}{\psi_{n}^{(1)}}=\sum_{k\neq n}\dfrac{\bra{\psi_{k}^{(0)}}\hat{H}^{(1)}\ket{\psi_{n}^{(0)}}}{E_n^{(0)}-E_k^{(0)}}\ket{\psi_k^{(0)}}\]
这一求和隐含的条件是$\psi_{n}^{(0)}$与$\psi_{n}^{(1)}$的正交性.由于$\psi_n$是归一的,因此
\[\inprod{\psi_{n}^{(0)}+\lambda\psi_{n}^{(1)}+o(\lambda)}{\psi_{n}^{(0)}+\lambda\psi_{n}^{(1)}+o(\lambda)}=\inprod{\psi_n^{(0)}}{\psi_n^{(0)}}+\lambda\left(\inprod{\psi_n^{(1)}}{\psi_n^{(0)}}+\inprod{\psi_n^{(0)}}{\psi_n^{(1)}}\right)+o(\lambda)=1\]
上式对所有$\lambda$都成立,因此只能有
\[\inprod{\psi_n^{(1)}}{\psi_n^{(0)}}+\inprod{\psi_n^{(0)}}{\psi_n^{(1)}}=0\]
在假定波函数是实值函数的前提下可得$\inprod{\psi_n^{(0)}}{\psi_n^{(1)}}=0$,即二者正交.
\subsubsection{含时微扰}
含时薛定谔方程为
\[\i\hbar\p_t\iPsi(\vec{r},t)=\hat{H}\iPsi(\vec{r},t)\]
如果$\hat{H}$不含时,并且$\psi(\vec{r})$是$\hat{H}$的定态薛定谔方程
\[\hat{H}\psi(\vec{r})=E\psi(\vec{r})\]
的解,那么含时波函数即为
\[\iPsi(\vec{r},t)=\psi(\vec{r})\exp\left(-\dfrac{\i Et}{\hbar}\right)\]
\indent 现在,将$\hat{H}$的含时部分视作微扰:
\[\hat{H}=\hat{H}^{(0)}+\hat{H}^{(1)}(t)\]
在\tbf{偶极近似}下,假定原子处在匀强电磁场中,那么微扰项
\[\hat{H}^{(1)}(t)=-\hat{\mu}_zE\cos\omega t\]
其中$z$方向指向电磁场的振荡方向, $\hat{\mu}_z$为偶极算符.\\
\indent 考虑一阶微扰波函数,与定态微扰相似,可以写作
\[\ket{\iPsi_n^{(1)}(t)}=\sum_{k\neq n}a_k(t)\ket{\iPsi^{(0)}_k(t)}\]
类比定态微扰(实际上推导过程是一样的)可得
\[a_k(t)=\dfrac{1}{\i\hbar}\int_0^{t}\bra{\iPsi_k^{(0)}(\tau)}\hat{H}^{(1)}(\tau)\ket{\iPsi_n^{(0)}(\tau)}\d\tau\]
将$\iPsi(\vec{r},t)=\psi_k(\vec{r})\exp\left(-\dfrac{\i E_k^{(0)}t}{\hbar}\right)$代入可得
\[a_k(t)=\dfrac{1}{\i\hbar}\int_0^{t}\bra{\psi_k^{(0)}}\hat{H}^{(1)}(\tau)\ket{\psi_n^{(0)}}\e^{\i\omega_{kn}\tau}\d\tau,\quad\omega_{kn}=\dfrac{E_k^{(0)}-E_n^{(0)}}{\hbar}\]
\indent 在光谱学中,上式即为\tbf{跃迁概率辐}.跃迁概率即为
\[P_k(t)=\left|a_k(t)\right|^2\]
定义\tbf{跃迁速率}
\[w_k(t)=\dfrac{\d P_k}{\d t}\]
在光致跃迁中
\[\hat{H}^{(1)}(t)=-\hat{\mu}_zE\cos\omega t\]
于是
\[a_k(t)=\mathcal{E}\bra{\psi_k^{(0)}}\hat{\mu}_z\ket{\psi_n^{(0)}}\dfrac{1}{\i\hbar}\int_0^{t}\e^{\i\omega_{kn}\tau}\cos\omega\tau\d\tau\]
从而
\[w_k(t)\propto\mathcal{E}^2\left|\bra{\psi_k^{(0)}}\hat{\mu}_z\ket{\psi_n^{(0)}}\right|^2\]
于是跃迁速率(正比于吸收/发射强度)正比于偶极矩矩阵元的模长的平方.定义\tbf{跃迁偶极矩}
\[\mu_{kn}=\bra{\psi_k^{(0)}}\hat{\mu}_z\ket{\psi_n^{(0)}}\]
\subsection{偶极近似}
\subsubsection{泡利方程}
电场中的薛定谔方程,即泡利方程为
\[\i\hbar\dfrac{\p}{\p t}\psi=\hat{H}\psi,\quad\hat{H}=\dfrac{1}{2m}\left[\symbf\sigma\cdot\left(\hat{\vec{p}}-e\vec{A}\right)\right]^2+e\phi\]
其中$\symbf\sigma$是泡利矩阵, $\vec{A}$是磁矢势.根据泡利矩阵的性质有
\[(\symbf\sigma\cdot\vec{a})(\symbf\sigma\cdot\vec{b})=\vec{a}\cdot\vec{b}+\i\symbf\sigma\cdot(\vec{a}\times\vec{b})\]
于是
\[\left[\symbf\sigma\cdot\left(\hat{\vec{p}}-e\vec{A}\right)\right]^2=\left(\vec{p}-e\vec{A}\right)^2+\i\sigma\left(\vec{p}-e\vec{A}\right)\times\left(\vec{p}-e\vec{A}\right)\]
而
\[\left(\vec{p}-e\vec{A}\right)\times\left(\vec{p}-e\vec{A}\right)\psi=-e\hat{\vec{p}}\times(\vec{A}\psi)-e\vec{A}\times(\hat{\vec{p}}\psi)=\i e\hbar\left[\nabla\times(\vec{A}\psi)+\vec{A}\times(\nabla\psi)\right]=\i e\hbar(\nabla\times\vec{A})\psi\]
从而泡利方程可以改写为
\[\hat{H}=\dfrac{1}{2m}\left(\hat{\vec{p}}-e\vec{A}\right)^2+e\phi-\dfrac{e\hbar}{2m}\symbf\sigma\cdot\mat{B},\quad\mat{B}=\nabla\times\vec{A}\]
根据磁矢势与磁场的关系可知上式中的$\mat{B}$即为磁场.
\subsubsection{偶极近似}
采用处理多电子原子相同的办法,考虑电磁场中的多电子原子的哈密顿量(在原子单位制下):
\[\hat{H}(t)=\dfrac12\sum_{j=1}^{N}\left[\dfrac{\nabla_j}{\i}+\vec{A}(\vec{r}_j,t)\right]^2+\sum_{j=1}^{N}v_0(\vec{r}_j)+\mu_B\sum_{j=1}^{N}\symbf\sigma\cdot\nabla_J\times\vec{A}(\vec{r}_j,t)+\dfrac12\sum_{j\neq k}^{N}\dfrac{1}{\left|\vec{r}_j-\vec{r}_k\right|}\]
在这里我们做了如下近似
\begin{enumerate}[label=(\alph*)]
  	\item 电磁场被经典处理,要求光子数目众多.
  	\item 非相对论处理,要求电子速度远低于光速.
\end{enumerate}
此外,由于通常电磁波的波长远大于原子的尺寸,因此近似地认为$\vec{A}$在局部空间内为常数,不依赖于$\vec{r}_j$.在此近似下,泡利方程为
\[\i\hbar\dfrac{\p}{\p t}\iPsi=\left(\dfrac12\sum_{j=1}^{N}\left[\dfrac{\nabla_j}{\i}+\vec{A}(t)\right]^2+\sum_{j=1}^{N}v_0(\vec{r}_j)+\dfrac12\sum_{j\neq k}^{N}\dfrac{1}{\left|\vec{r}_j-\vec{r}_k\right|}\right)\iPsi\]
做酉变换
\[\iPsi'(t)=\e^{\i\phi(t)}\iPsi(t),\quad\phi(t)=\sum_{j=1}^{N}\vec{r}_j\cdot\vec{A}(t)\]
这样原方程变为
\[\i\hbar\dfrac{\p}{\p t}\iPsi'=\hat{H}'(t)\iPsi',\quad\hat{H}'(t)=\e^{\i\phi(t)}\hat{H}(t)\e^{-\i\phi(t)}-\dfrac{\p\phi(t)}{\p t}\]
整理可得
\[\i\dfrac{\p}{\p t}\iPsi'=\left(-\sum_{j=1}^{N}\dfrac{\nabla_j^2}{2}+\sum_{j=1}^{N}\left[\vec{r}_j\cdot\vec{E}(t)+v_0(\vec{r}_j)\right]+\dfrac{1}{2}\sum_{j\neq k}\dfrac{1}{\left|\vec{r}_j-\vec{r}_k\right|}\right)\iPsi',\quad\vec{E}(t)=-\dfrac{\p\vec{A}(t)}{\p t}\]
\subsection{原子光谱}
\subsubsection{类氢原子的光谱}
我们已经知道跃迁的强度正比于跃迁偶极矩.而
\[\mu_{z,fi}=-e\int\iPsi_f^\ast z\iPsi_i\d\tau\]
代入类氢原子轨道可得
\[\mu_{z,fi}=-e\int_0^{\infty}\int_0^{2\pi}\int_0^\pi R_{n_f,l_f}Y_{l_f,m_{l,f}}^\ast\sqrt{\dfrac{4\pi}{3}}rY_{1,0}R_{n_i,l_i}Y_{l_i,m_{l,i}}r^2\d r\sin\theta\d\theta\d\phi\]
整理可得
\[\mu_{z,fi}=-\sqrt{\dfrac{4\pi}{3}}e\int_0^{\infty}R_{n_f,l_f}r^3R_{n_i,l_i}\d r\int_0^{2\pi}\int_0^\pi Y_{l_f,m_{l,f}}^\ast Y_{1,0}Y^{l_i,m_{l,i}}\sin\theta\d\theta\d\phi\]
定义一般的角向积分
\[I=\int_0^{2\pi}\int_0^\pi Y_{l_f,m_{l,f}}^\ast Y_{l,m}Y^{l_i,m_{l,i}}\sin\theta\d\theta\d\phi\]
$I\neq0$当且仅当$l_f=l_i\pm l$, $m_{l,f}=m_{l,i}\pm m$.对于类氢原子有$l=1$, $m=0$,因此
\begin{theorem}[类氢原子的跃迁选律]
	类氢原子的跃迁满足
	\[\Delta l=\pm1,\quad\Delta m_l=0,\pm1\]
\end{theorem}
\subsubsection{氦原子的光谱}
氦原子的激发态的电子不在同一个轨道上(激发两个电子所需的能量高于电离能,因此至多只有一个电子被激发,并且根据选律可知电子组态为$1\text{s}^1n\text{l}^1$),因此电子自旋的排布方式意味着氦原子的同一电子组态有单重态和三重态,并且一般而言三重态的能量低于单重态.\tbf{单重态和三重态之间也不能发生跃迁}.
\subsubsection{旋轨耦合}
电子绕核旋转产生轨道角动量,自旋产生自旋角动量.两者产生的磁矩发生相互作用,倾向于方向相反,这一相互作用称作\tbf{旋轨耦合(Spin-Orbit Coupling, SOC)}.\\
\indent 定义总角动量为轨道角动量与自旋角动量的矢量和:
\[\vec{j}=\vec{l}+\vec{s}\]
在经典力学中,磁矩在磁场中的能量为
\[E=-\symbf\mu\cdot\vec{B}\]
而
\[\symbf\mu=\dfrac{e}{2m}\vec{s},\quad\mat{B}=\dfrac{\mu_0 q}{8\pi m}\cdot\dfrac{3(\vec{l}\cdot\vec{r})\vec{r}-\vec{l}}{r^3}\propto\vec{l}\]
从而
\[E=-\symbf\mu\cdot\vec{B}\propto\vec{s}\cdot\vec{l}\]
\indent 根据$\vec{j}$的定义有
\[\vec{j}^2=\vec{l}^2+2\vec{s}\cdot\vec{l}+\vec{s}^2\]
于是
\[\vec{s}\cdot\vec{l}=\dfrac12\left(\vec{j}^2-\vec{l}^2-\vec{s}^2\right)\]
根据量子力学中角动量和对应量子数的关系可得
\[\vec{s}\cdot\vec{l}=\dfrac12\left[j(j+1)-l(l+1)-s(s+1)\right]\hbar^2\]
其中$l$, $s$和$j$分别为轨道角量子数,自旋磁量子数和总角动量量子数.最终可得旋轨耦合能为
\[E_{l,s,j}=\dfrac12hc\tilde{A}\left[j(j+1)-l(l+1)-s(s+1)\right]\]
其中$\tilde{A}$是\tbf{旋轨耦合常数},与波数具有相同的量纲.旋轨耦合常数的大小正比于核电荷数的四次方,因此重原子核的旋轨耦合不可以忽略.\\
\indent 旋轨耦合的将导致同一谱项的光谱支项发生能级分裂.钠光灯光谱的双线结构就是由$^2\text{S}_{1/2}$跃迁至$^2\text{P}_{1/2}$和$^3\text{P}_{3/2}$所形成的两条谱线.
\subsection{光谱项和光谱支项的符号}
\tbf{光谱支项}的记号如下:
\[^{2S+1}L_J\]
其中$2S+1$为自旋多重度, $S$即总自旋量子数. $L$为总轨道角动量量子数对应的符号($L=0,1,2,\cdots$分别对应于S, P, D等), $J$为总角动量量子数.如果不考虑旋轨耦合作用即为\tbf{光谱项},其记号为
\[^{2S+1}L\]
\indent 总轨道角动量量子数$L$通过各个电子的角动量耦合得到.对于两个电子的情形, $L$的取值通过\tbf{Clebsch-Gordan级数}得到:
\[L=l_1+l_2,l_1+l_2-1,\cdots,\left|l_1-l_2\right|\]
其中$l_1$和$l_2$为两个电子的角动量量子数.如果有多个电子,依次按上面的步骤耦合即可.
\indent 总自旋量子数$S$通过各个电子的自旋角动量耦合得到,其耦合方法与角动量的耦合类似.需要注意的是,在完成上述操作之后,不能出现违反泡利不相容原理的电子.\\
\indent 多电子原子的角动量有两种耦合方式:
\begin{enumerate}[label=\tbf{(\alph*)}]
	\item 总轨道角动量$\vec{L}$与总自旋角动量$\vec{S}$耦合(Russell-Saunders Coupling, RS Coupling).适用于旋轨耦合较弱的情况(轻核).总角动量量子数的取值为:
	\[J=L+S,L+S-1,\cdots,|L-S|\]
	\item 每个电子的轨道角动量$\vec{s}$和自旋角动量$l$耦合成$\vec{j}$,然后不同电子再耦合成总角动量$\vec{J}$(jj-Coupling).适用于适用于旋轨耦合较强的情况(重核).
\end{enumerate}
利用光谱项的记号,洪特规则可以重写如下:
\begin{theorem}[洪特规则]
	\begin{enumerate}[label=\tbf{\alph*.}]
		\item 给定构型下,多重度最大的光谱项能量最低.
		\item 给定$S$, $L$最大的光谱项能量最低.
		\item 小于半满, $J$最小能量最低;大于半满, $J$最大能量最低.
	\end{enumerate}
\end{theorem}
在同一光谱项对应的所有光谱支项中, 少于半满的情况下$J$越小能量越低,而多于半满的情况下$J$越大能量越低.

\end{document}